% SuperClase_Econometria_InflacionChile_COMPLETO.tex
\documentclass[aspectratio=169,10pt]{beamer}

% ----------- Idioma y fuentes -----------
\usepackage[spanish,es-tabla]{babel}
\usepackage[T1]{fontenc}
\usepackage[utf8]{inputenc}
\usepackage{inconsolata}
\usepackage{lmodern}

% ----------- Colores UdeC -----------
\usepackage[table]{xcolor}
\definecolor{UdeCBlue}{HTML}{003366}
\definecolor{UdeCYellow}{HTML}{FDB813}
\definecolor{LightGray}{HTML}{F5F5F5}
\definecolor{CodeBG}{HTML}{F2F4F7}
\definecolor{CodeFrame}{HTML}{D0D5DD}

% ----------- Tema Beamer -----------
\usetheme{Madrid}
\usecolortheme{default}

\setbeamercolor{title}{fg=UdeCBlue}
\setbeamercolor{frametitle}{fg=UdeCBlue,bg=UdeCBlue!6}
\setbeamercolor{structure}{fg=UdeCBlue}
\setbeamercolor{block title}{fg=UdeCBlue,bg=UdeCBlue!10}
\setbeamercolor{block body}{fg=black,bg=UdeCBlue!2}
\setbeamercolor{alerted text}{fg=UdeCYellow!80!black}

\setbeamertemplate{itemize items}[circle]
\setbeamertemplate{enumerate items}[circle]
\setbeamertemplate{navigation symbols}{}

% ----------- Encabezado / pie -----------
\setbeamercolor{section in head/foot}{fg=white,bg=UdeCBlue}
\setbeamercolor{author in head/foot}{fg=white,bg=UdeCBlue!80}
\setbeamercolor{title in head/foot}{fg=white,bg=UdeCBlue}
\setbeamercolor{date in head/foot}{fg=white,bg=UdeCBlue!80}

\setbeamertemplate{footline}{%
  \leavevmode%
  \hbox{%
    \begin{beamercolorbox}[wd=.33\paperwidth,ht=2.5ex,dp=1ex,left]{author in head/foot}%
      \hspace*{1.5ex}\insertshortauthor
    \end{beamercolorbox}%
    \begin{beamercolorbox}[wd=.44\paperwidth,ht=2.5ex,dp=1ex,center]{title in head/foot}%
      \textbf{MACROECONOMÍA (580323)}
    \end{beamercolorbox}%
    \begin{beamercolorbox}[wd=.23\paperwidth,ht=2.5ex,dp=1ex,right]{date in head/foot}%
      \insertframenumber{} / \inserttotalframenumber\hspace*{1.5ex}
    \end{beamercolorbox}%
  }%
  \vskip0pt%
}

\setbeamertemplate{headline}{%
  \begin{beamercolorbox}[wd=\paperwidth,ht=1.8ex,dp=0.8ex,left]{section in head/foot}
    \hspace*{1.2ex}\insertsectionhead
  \end{beamercolorbox}
}

% ----------- Hyperref -----------
\usepackage{hyperref}
\hypersetup{colorlinks=true, linkcolor=UdeCBlue, urlcolor=UdeCBlue}

% ----------- Código -----------
\usepackage{listings}
\usepackage{listingsutf8}
\lstdefinelanguage{Python}{
  keywords={from, import, as, if, elif, else, for, while, in, try, except, finally, with, return, def, class, pass, break, continue, lambda, True, False, None, and, or, not},
  sensitive=true,
  comment=[l]\#,
  morecomment=[s]{\"\"\"}{\"\"\"},
  morestring=[b]',
  morestring=[b]"
}
\lstset{
  language=Python,
  inputencoding=utf8,
  basicstyle=\ttfamily\footnotesize,
  backgroundcolor=\color{CodeBG},
  frame=single,
  rulecolor=\color{CodeFrame},
  frameround=tttt,
  keywordstyle=\color{UdeCBlue}\bfseries,
  commentstyle=\color{gray!70}\itshape,
  stringstyle=\color{UdeCYellow!70!black},
  showstringspaces=false,
  numbers=left,
  numberstyle=\tiny\color{gray},
  numbersep=6pt,
  columns=fullflexible,
  keepspaces=true,
  aboveskip=0.8em,
  belowskip=0.8em
}

% ----------- Datos de portada -----------
\title[Super Clase Econometría]{Super Clase de Econometría Aplicada\\ Predicción de Inflación en Chile}
\author[]{Departamento de Ingeniería Industrial\\ \textbf{Macroeconomía (580323)}}
\institute[]{Universidad de Concepción}
\date{\today}

% --------------------------------------- DOCUMENTO ---------------------------------------
\begin{document}

% ========================== PORTADA ==========================
\begin{frame}[plain]

  % ----- TÍTULO COMPLETO (oficial + tu título personal) -----
  \begin{center}
    {\Large \bfseries Ayudantía para Trabajo Semestral}\\[0.3em]
    {\large Predicción de Inflación en Chile}\\[1.2em]

    {\Large \color{UdeCBlue}\textbf{Introducción a la Econometría}}\\[0.4em]
    {\small Por Ignacio Garrido U.}\\[2em]
  \end{center}

  % ----- LÍNEA Y OBJETIVO -----
  \vfill
  \begin{center}
    \color{UdeCBlue}\rule{0.65\linewidth}{1pt}\\[0.4em]
    \small 
    \textbf{Objetivo}: Construir, interpretar y validar un modelo para predecir inflación mensual (IPC) 
    con foco en \alert{teoría}, \alert{datos} y \alert{validación fuera de muestra}.
  \end{center}

\end{frame}



% =====================================================================
% ========================== SECCIÓN 0: INTRODUCCIÓN ==================
% =====================================================================

\section{Introducción general}

\begin{frame}{¿Qué es Econometría?}
  \begin{itemize}
    \item Integración de \textbf{Economía + Estadística + Matemáticas + Programación}.
    \item Responde preguntas empíricas:
      \begin{itemize}
        \item ¿Cómo se relacionan las variables?
        \item ¿Podemos predecir eventos económicos?
      \end{itemize}
    \item La tecnología (Python, R, Stata) permite:
      \begin{itemize}
        \item Modelos complejos,
        \item Bases de datos masivas,
        \item Automatizar validaciones.
      \end{itemize}
  \end{itemize}
\end{frame}

\begin{frame}{Econometría en macro y microeconomía}
  \begin{itemize}
    \item \textbf{Micro}: hogares, empresas, consumo, deuda.
    \item \textbf{Macro}: inflación, empleo, PIB, tipo de cambio.
    \item Dos fines:
      \begin{itemize}
        \item \textbf{Explicar} relaciones.
        \item \textbf{Predecir} futuro.
      \end{itemize}
    \item Pilar clave: teoría + datos + validación.
  \end{itemize}
\end{frame}

% =====================================================================
% ======================= SECCIÓN 1B: LITERATURA ======================
% =====================================================================

\section{Lectura y búsqueda de literatura}

\begin{frame}{Dónde buscar literatura confiable}
  \begin{itemize}
    \item Evitar blogs o documentos sin revisión.
    \item Portales recomendados:
      \begin{itemize}
        \item \textbf{Web of Science (WOS)}
        \item \textbf{SCOPUS}
        \item \textbf{SciELO}
        \item \textbf{IDEAS/RePEc}
      \end{itemize}
    \item Google Scholar sirve, pero requiere criterio.
  \end{itemize}
\end{frame}

\begin{frame}{Cómo leer un paper correctamente}
  \begin{itemize}
    \item No basta con mencionarlo.
    \item Lectura debe ser \textbf{cronológica y conectada}.
    \item Orden recomendado:
      \begin{enumerate}
        \item Problema económico.
        \item Teoría que lo respalda.
        \item Datos: fuente, frecuencia.
        \item Modelo aplicado.
        \item Resultados.
        \item Limitaciones.
      \end{enumerate}
  \end{itemize}
\end{frame}

% =====================================================================
% ========================== SECCIÓN 2: DATOS =========================
% =====================================================================

\section{Datos: fuentes y preparación}

\begin{frame}{Fuentes de datos recomendadas}
  \begin{itemize}
    \item \textbf{Banco Central de Chile}: TPM, dólar.
    \item \textbf{INE}: IPC mensual.
    \item \textbf{SII}: series tributarias.
    \item \textbf{CMF}: crédito, bancos.
    \item Páginas útiles:
      \begin{itemize}
        \item Datosmacro.com (no oficial).
      \end{itemize}
  \end{itemize}
\end{frame}

\begin{frame}{Frecuencia y consistencia}
  \begin{itemize}
    \item El IPC es \textbf{mensual}.
    \item Variables explicativas deben:
      \begin{itemize}
        \item Estar alineadas a la misma frecuencia.
        \item No tener fechas faltantes.
      \end{itemize}
    \item Unificar unidades y escalas antes de modelar.
  \end{itemize}
\end{frame}

% =====================================================================
% ========================== SECCIÓN 3: DATASET =======================
% =====================================================================

\section{Construcción del dataset}

\begin{frame}{Variable dependiente y explicativas}
  \begin{block}{Dependiente}
    IPC mensual (\% m/m).
  \end{block}
  \begin{block}{Explicativas}
    \begin{itemize}
      \item TPM rezagada,
      \item USD/CLP rezagado,
      \item Cobre,
      \item Rezagos del IPC,
      \item Expectativas económicas.
    \end{itemize}
  \end{block}
  \begin{itemize}
    \item Clave: misma frecuencia y misma cantidad de datos.
  \end{itemize}
\end{frame}

\begin{frame}[fragile]{Formato de datos}
\begin{itemize}
  \item CSV (recomendado): liviano y universal.
  \item XLSX: más pesado, depende de Excel.
  \item Arrays/listas: eficientes pero menos descriptivos.
\end{itemize}
\end{frame}

\begin{frame}[fragile]{Carga y creación de rezagos}
\begin{lstlisting}
df = pd.read_csv("dataset_inflacion.csv", parse_dates=["fecha"])
df = df.sort_values("fecha").reset_index(drop=True)

for col in ["ipc_mom","tpm","usdclp"]:
    df[f"{col}_lag1"] = df[col].shift(1)
    df[f"{col}_lag2"] = df[col].shift(2)

df = df.dropna().reset_index(drop=True)
df.head()
\end{lstlisting}
\end{frame}

% =====================================================================
% ========================== SECCIÓN 4: MODELOS =======================
% =====================================================================

\section{Modelos y lectura de resultados}

\begin{frame}[fragile]{Estimación OLS (modelo base)}
\begin{lstlisting}
X = df[["ipc_mom_lag1","tpm_lag1","usdclp_lag1"]]
X = sm.add_constant(X)
y = df["ipc_mom"]

ols = sm.OLS(y, X).fit(cov_type="HC1")
print(ols.summary())
\end{lstlisting}
\end{frame}

\begin{frame}{Interpretación del resumen}
  \begin{itemize}
    \item \textbf{R²}: proporción explicada.
    \item \textbf{R² ajustado}: corrige sobreparametrización.
    \item \textbf{Prueba F}: significancia global.
    \item \textbf{Coeficientes}: dirección y magnitud.
    \item \textbf{p-valores}: evidencia estadística.
    \item \textbf{Errores robustos HC1}: corrigen heterocedasticidad.
  \end{itemize}
\end{frame}

\begin{frame}{Ejemplo de resultados}

\centering
\small
\begin{tabular}{lccc}
\hline
 & Coeficiente & Error Std. & p-valor \\
\hline
\textbf{Constante}       & 0.152 & 0.041 & 0.001 \\
\textbf{IPC$_{t-1}$}     & 0.473 & 0.052 & 0.000 \\
\textbf{TPM$_{t-1}$}     & 0.018 & 0.007 & 0.014 \\
\textbf{USD/CLP$_{t-1}$} & 0.00042 & 0.00019 & 0.028 \\
\hline
\textbf{R$^2$} & \multicolumn{3}{c}{0.52} \\
\textbf{R$^2$ ajustado} & \multicolumn{3}{c}{0.51} \\
\textbf{F-statistic} & \multicolumn{3}{c}{34.21 (p = 0.000)} \\
\textbf{N} & \multicolumn{3}{c}{108} \\
\hline
\end{tabular}

\vspace{0.4cm}
\small Modelo estimado con errores estándar robustos (HC1).

\end{frame}


% =====================================================================
% ========================== SECCIÓN 5: ROBUSTEZ ======================
% =====================================================================

\section{Diagnóstico y robustez}

\begin{frame}[fragile]{Normalidad y autocorrelación: ¿qué mide cada prueba?}

\begin{itemize}
  \item \textbf{Jarque–Bera (JB)}
    \begin{itemize}
      \item Evalúa si los residuos tienen \textbf{asimetría} y \textbf{curtosis} compatibles con una distribución normal.
      \item \textbf{H0}: los residuos son normales.
      \item \textbf{Interpretación}: p-valor bajo $\Rightarrow$ evidencia contra normalidad.
    \end{itemize}

  \item \textbf{Anderson–Darling (AD)}
    \begin{itemize}
      \item Prueba de normalidad enfocada especialmente en las \textbf{colas} de la distribución.
      \item \textbf{H0}: residuos normales.
      \item Estadístico alto $\Rightarrow$ residuos no normales.
    \end{itemize}

  \item \textbf{Durbin–Watson (DW)}
    \begin{itemize}
      \item Detecta \textbf{autocorrelación de primer orden} en los residuos.
      \item Valores típicos:
        \begin{itemize}
          \item 2 $\Rightarrow$ sin autocorrelación
          \item $< 2$ $\Rightarrow$ autocorrelación positiva
          \item $> 2$ $\Rightarrow$ autocorrelación negativa
        \end{itemize}
    \end{itemize}
\end{itemize}

\end{frame}

\begin{frame}[fragile]{Normalidad y autocorrelación: ¿qué mide cada prueba?}


\begin{lstlisting}
resid = ols.resid
print("JB:", jarque_bera(resid))
print("AD:", normal_ad(resid))
print("DW:", durbin_watson(resid))
\end{lstlisting}

\end{frame}
\begin{frame}[fragile]{Heterocedasticidad, VIF y RESET: significado e interpretación}

\begin{itemize}

  \item \textbf{Breusch–Pagan (BP)}
    \begin{itemize}
      \item Detecta si la varianza de los residuos es \textbf{no constante} (heterocedasticidad).
      \item \textbf{H0}: homocedasticidad.
      \item p-valor bajo $\Rightarrow$ heterocedasticidad presente.
      \item Solución frecuente: usar errores robustos (HC0, HC1).
    \end{itemize}

  \item \textbf{VIF (Variance Inflation Factor)}
    \begin{itemize}
      \item Mide \textbf{multicolinealidad} entre regresores.
      \item Valores típicos:
        \begin{itemize}
          \item VIF = 1 $\Rightarrow$ sin colinealidad.
          \item 5–10 $\Rightarrow$ colinealidad moderada.
          \item $> 10$ $\Rightarrow$ colinealidad severa.
        \end{itemize}
      \item Colinealidad elevada $\Rightarrow$ coeficientes inestables y errores grandes.
    \end{itemize}

  \item \textbf{RESET de Ramsey} (si decides mencionarlo)
    \begin{itemize}
      \item Evalúa \textbf{mala especificación funcional}.
      \item Detecta si faltan términos no lineales o variables omitidas.
    \end{itemize}

\end{itemize}

\end{frame}

\begin{frame}[fragile]{Heterocedasticidad, VIF y RESET: significado e interpretación}


\begin{lstlisting}
Xtr = sm.add_constant(df[["ipc_mom_lag1","tpm_lag1","usdclp_lag1"]])
mod = sm.OLS(df["ipc_mom"], Xtr).fit()

print("BP:", het_breuschpagan(mod.resid, Xtr))

[(c, variance_inflation_factor(Xtr.values,i))
 for i,c in enumerate(Xtr.columns)]
\end{lstlisting}

\end{frame}
% =====================================================================
% ========================== SECCIÓN 6: VALIDACIÓN ====================
% =====================================================================

\section{Validación fuera de muestra}

\begin{frame}{Predicción del último período}
  \begin{itemize}
    \item Entrenar con toda la base menos el último mes.
    \item Predecir el último mes conocido.
    \item Si predicción $\approx$ real $\Rightarrow$ modelo estable.
    \item Si no $\Rightarrow$ revisar variables, rezagos o forma funcional.
  \end{itemize}
\end{frame}

% =====================================================================
% ========================== SECCIÓN 7: CONSEJOS ======================
% =====================================================================

\section{Recomendaciones finales}

\begin{frame}{Uso responsable de IA}
  \begin{itemize}
    \item Útil para:
      \begin{itemize}
        \item Explicar errores de código,
        \item Revisar sintaxis,
        \item Entender conceptos.
      \end{itemize}
    \item No confiable en:
      \begin{itemize}
        \item Referencias y papers (puede inventarlos),
        \item Datos oficiales,
        \item Cálculos econométricos avanzados.
      \end{itemize}
    \item Úsenla como apoyo, no como verdad absoluta.
  \end{itemize}
\end{frame}

\begin{frame}{Cierre}
  \begin{itemize}
    \item Mejoren, comparen, validen.
    \item Pregunten con tiempo si tienen dudas.
    \item Este trabajo los inicia en el mundo real de los datos.
  \end{itemize}
  \vfill
  \begin{center}
    {\Large \textbf{\color{UdeCBlue}y... Bienvenidos al mundo de los datos}}
  \end{center}
\end{frame}

% ------------------ FIN DEL DOCUMENTO ------------------
\end{document}
